\section{Introduction}\label{sec:section_1}

Arc spaces describe formal power series solutions (in one variable) to 
polynomial equations. They first appeared in the work of Nash (published later as \cite{nash}), who investigated their relation to some intrinsic data of a resolution of singularities of a fixed algebraic variety. He asked whether there is a bijection between the irreducible components of the arc space based at the singular locus and the set of essential divisors. While this so-called `Nash problem' is still actively studied (see for instance the recent papers \cite{monique_ana,plenat_spivakovsky,pe,fernandez} and the overview \cite{ishii_intro}), in the last decade arc spaces have gained much interest from algebraic geometers, through their role in motivic integration and their utility in birational geometry.
Arc spaces show strong relations with combinatorics as well. In the present text we
indicate how to exploit this connection both for geometric as well as
combinatorial benefit. In particular, we expose a surprising connection with the well-known
Rogers-Ramanujan identities.\\

Let us emphasize the main algebraic and combinatorial aspects
presented here. First, we suggest to study local algebro-geometric
properties of algebraic (or analytic) varieties via natural
Hilbert-Poincar\'{e} series attached to arc spaces. In contrast to
already existing such series this one is sensitive to the
non-reduced structure of the arc space. Second, we propose
to derive identities between partitions by looking at suitable ideals
in a polynomial ring in countably many variables endowed with a
natural grading. Connecting both ideas will demand handling
Gr\"{o}bner basis in countably many variables, a problem which has
been successfully dealt with in different contexts over the last
years (see \cite{hillar_sullivant,draisma}). In the present situation -- that is for very specific ideals --
salvation from the natural obstruction of being infinitely generated
comes in the shape of a derivation making the respective ideals
differential.\\

We briefly indicate the connection between arc spaces and
partitions. Let $f\in k[x_1,\ldots, x_n]$ be a polynomial in $n$ variables $x_1,\ldots, x_n$
with coefficients in a field $k$. We denote the formal power series ring in one
variable $t$ over the field $k$ by $k[[t]]$. The \emph{arc
 space} $X_\infty$ of the algebraic variety $X$ defined by $f$ is the set of
power series solutions \begin{math} x(t)=(x_1(t),\ldots, x_n(t))\in k[[t]]^n
\end{math} to the equation $f(x(t))=0$. This set
turns out to be eventually algebraic in the sense that it is given by
polynomial equations (though there are countably many of
them). Indeed, expanding $f(x(t))$ as a power series in $t$
gives 
\begin{displaymath}
  f(x(t)) = F_0 + F_1t + F_2t^2 + \cdots
\end{displaymath}
where the $F_i$ are polynomials in the coefficients of $t$ in $x(t)$. Therefore, a given vector of formal
power series $a(t)\in k[[t]]^n$ is an element of the arc space
$X_\infty$ if and only if its coefficients fulfill
the equations $F_0,F_1,\dots$. Algebraically the corresponding set of
solutions is described by its \emph{coordinate algebra}
\begin{displaymath}
  J_\infty(X)= k[x_j^{(i)}; 1\leq j \leq n, i\in \N]/(F_0,F_1,\ldots),
\end{displaymath}
where $\mathbb{N}=\{0,1,2,\ldots \}$. The variable $x_j^{(i)}$ corresponds to the coefficient of $t^i$
in $x_j(t)$. We will mostly be interested in the case where $a(0)$
is a point on $X$ (without loss of generality we may assume that this is the origin). The resulting algebra, obtained from $J_\infty(X)$
by substituting $x_j^{(0)}=0$, is called the \emph{focussed
  arc algebra} and denoted by $J_\infty^0(X)$; we write $f_i$ for
the image of $F_i$ under this substitution:
\begin{displaymath}
  J_\infty^0(X)=k[x_j^{(i)}; 1\leq j \leq n, i\geq 1]/(f_1,f_2,\ldots).
\end{displaymath}
This algebra is naturally graded by the
weight function $\wt\, x_j^{(i)}=i$ since $f_{\ell}$ is homogeneous of weight $\ell$. In
the special case of $n=1$ we will write $y_i$ instead of
$x_1^{(i)}$. Integer partitions arise naturally when computing weights
of monomials in $J^0_\infty(\A^1)$. Recall that a \emph{partition} of
$m\in \N$ is an $r$-tuple of positive integers
$\lambda_1 \leq \lambda_2\leq \cdots \leq \lambda_r$ with $\lambda_1 +
\cdots + \lambda_r =m$. The $\lambda_i$ are the \emph{parts of the
  partition} and $r$ is its \emph{length}. A monomial
$y_1^{\alpha_1}\cdots y_e^{\alpha_e}$ has
weight $\alpha_1\cdot 1 + \cdots + \alpha_e\cdot
e$. Asking for the number of monomials (up to coefficients) of some
weight $m$ is thus asking for the number of partitions of $m$. This is 
precisely what we capture when computing the Hilbert-Poincar\'{e} series of
$J_\infty^0(\A^1)$. In general, the Hilbert-Poincar\'{e} series of $J_{\infty}^0(X)$ is defined as
\begin{displaymath}
  \HP_{J_\infty^0(X)}(t) = \sum_{j=0}^\infty \dim_k
  \left(J_\infty^0(X) \right)_j\cdot t^j,
\end{displaymath}
where $\left(J_\infty^0(X) \right)_j$ denotes the $j$th homogeneous
component of $J_\infty^0(X)$. In the simple case of $X=\A^1$ we may use the generating function for partitions to represent
$\HP_{J_\infty^0(\A^1)}(t)$ by 
\begin{displaymath}
  \mathbb{H} := \prod_{i\geq 1} \frac{1}{1-t^i}.
\end{displaymath}  

By the general theory of Gr\"{o}bner basis $\HP_{J_\infty^0(X)}(t)$ is identical with the Hilbert-Poincar\'{e} series of the algebra 
\begin{displaymath}
  k[x_j^{(i)}; 1\leq j \leq n, i\geq 1]/L(I),
\end{displaymath}
where $L(I)$ denotes the \emph{leading ideal} of $I=(f_0,f_1,\ldots)$ (with respect to 
a chosen monomial ordering). The leading ideal is much simpler since
it is generated by monomials. Computing the Hilbert-Poincar\'{e}
series of the respective algebra corresponds to counting partitions, leaving out
those coming from weights of monomials in $L(I)$. In the simple
example of $f=y^2$, these will be all partitions without repeated or
consecutive parts. Such partitions are part of the well-known
\emph{Rogers-Ramanujan identity}: the number of partitions of $n$ into
parts congruent to 1 or 4 modulo 5 is equal to the number of
partitions of $n$ into parts that are neither repeated nor
consecutive (see \cite{andrews_partitions}). This gives:

\begin{Thm*}
Let $k$ be field of characteristic $0$. For $X\colon y^2=0$ we compute
\begin{displaymath}
 \HP_{J_{\infty}^{0}(X)}(t) = \prod_{\substack{i\geq 1 \\ i\equiv
    1,4 \bmod 5}} \frac{1}{1-t^i}.
\end{displaymath}
\end{Thm*}

More generally, we obtain using Gordon's generalizations of the
Rogers-Ramanujan identities for $X\colon y^n=0$, $n\geq 2$:

\begin{samepage}
\begin{Thm*}
\begin{displaymath}
\HP_{J_\infty^0(X)}(t)=\Hh\cdot
  \prod_{\substack{i\geq 1\\ i \equiv 0, n,n+1 \\ \bmod 2n+1}} (1-t^i).
\end{displaymath}
\end{Thm*}
\end{samepage}

The computation of $L(I)$ in these rather simple
looking cases is nontrivial. It is carried out in Section
\ref{sec:section_5}.\\


Moreover, standard techniques from commutative algebra allow to compute a recursion for the
Hilbert-Poincar\'{e} series in the case of $X\colon y^2=0$:

\begin{Prop*}
The generating series $\HP_{J_{\infty}^{0}(X)}(t)$ is the $t$-adic limit of the sequence of formal power series $A_d$ defined by:
$$A_1=A_2=1, \quad \text{and} \quad A_d = A_{d-1}+t^{d-2} A_{d-2}\text{ for } d\geq 3.$$
\end{Prop*}

The above proposition was first found in an empirical way in
\cite{andrews_baxter}, and it leads to the Rogers-Ramanujan
identities (see Section \ref{sec:section_5}).\\


Returning to the geometric aspect of the series attached to arc spaces, we
compute them in other interesting cases: for smooth
points, for rational double points of surfaces, and for normal
crossings singularities. In these cases the simple geometry of the jet
schemes permits to compute the Hilbert-Poincar\'{e} series defined
above, and we find the following (see
Propositions~\ref{prop:smooth_case}, \ref{cor:rational_double_point}
and \ref{prop:normal_crossings}).

\begin{Prop*} With the above introduced notation, we have:
\begin{itemize}
\item[(1)] If $\mathfrak{p}$ is a smooth point on a variety $X$ of dimension $d$, then 
\[ \HP_{J_{\infty}^{\mathfrak{p}}(X)} = \left(\prod_{i\geq 1} \frac{1}{1-t^i}\right)^d, \]
\item[(2)] if $X$ is a surface with a rational double point at $\mathfrak{p}$ then 
\[ \HP_{J_{\infty}^{\mathfrak{p}}(X)}(t) =  \left(\frac{1}{1-t}\right)^{3} \left(\prod_{i\geq 2} \frac{1}{1-t^i}\right)^2, \]
\item[(3)] if $X= \{ x_1\cdots x_{d+1} = 0\} \subset \mathbb{A}^{d+1}_k$, and $\mathfrak{p}$ is a point in the intersection of precisely $e$ components, then 
\[ \HP_{J_{\infty}^{\mathfrak{p}}(X)}(t) =  \left(\prod_{i=1}^{e-1} \frac{1}{1-t^i}\right)^{d+1} \left(\prod_{i\geq e} \frac{1}{1-t^i}\right)^d. \]
\end{itemize}
\end{Prop*}

In the past several generating series have been associated to
the arc space of a (singular) variety by Denef and Loeser (see
\cite{denef_loeser,denef_loeser_definable_sets}), in analogy with the
$p$-adic case. All those series are defined in the motivic setting and
take values in a power series ring with coefficients in (a
localization of) a Grothendieck ring. They are rational. Some of those
series encode more information about the singularities than
others. For a comparison between them see
\cite{nicaise_manuscripta,nicaise,helena_pedro}. While those series
are concerned with the reduced structure of the arc space, our series
is \emph{sensitive to the non-reduced structure} as well, although it is not
rational in general. We discuss more geometric motivation for introducing the above Hilbert-Poincar\'{e} series in Section~\ref{sec:section_3} (after Definition~\ref{def:arc_Hilbert}).\\

The structure of the paper is as follows: in Section~\ref{sec:section_2}, we define jet schemes and arc spaces and recall some basic facts about them.
In Section~\ref{sec:section_3} we introduce the arc
Hilbert-Poincar\'{e} series, we discuss its properties and we compute
it in particular cases. Section~\ref{sec:section_4} recalls basic
facts about partitions and the Rogers-Ramanujan
identities. Section~\ref{sec:section_5} is devoted to the main
theorem. For the convenience of the reader, we recall the facts about
Hilbert-Poincar\'{e} series and Gr\"obner bases that we use in the
paper in an appendix (Section~\ref{appendix}).

\subsection*{Acknowledgements} 
The authors thank the organizers of YMIS for preparing the ground for
this joint research. The first named author expresses his gratitude to
Herwig Hauser, Georg Regensburger and Josef Schicho for many useful
discussions on these topics; the second named author is grateful to
Monique Lejeune-Jalabert for introducing him to Hilbert-Poincar\'{e}
series; the third named author thanks Johannes Nicaise for explaining 
him several facts about arc spaces in detail and he thanks the
Institut des Hautes \'{E}tudes Scientifiques, where part of this work
was done, for hospitality.\\

